% GENERATED FILE. Edit canonical lessons, then run npm run generate:books.
% canonical-source-hash: fnv1a64:39c2dc7887756777
% canonical-lessons: SA-C226-grahah, SA-C226-dhumaketuh, SA-C226-pasuh, SA-C226-jantuh, SA-C226-paksi, SA-R226-first-pass-words-from-chapters-221-223, SA-R226-second-pass-words-from-chapters-224-226

\chapter{Planet, Comet, Animal, Creature, Winged creature}
\label{ch:sa-planet-comet-animal-creature-winged-creature}
\hlchaptermodality{226}

\begin{chapteropening}
\textbf{By the end of this chapter:} \emph{I can say a planet, a comet, an animal, a creature and a winged creature.}

Complete the last lesson of chapter 226: I can say a planet, a comet, an animal, a creature and a winged creature.
\end{chapteropening}

\section[\texorpdfstring{\sk{ग्रहः} (grahaḥ)}{ग्रहः (grahaḥ)}]{\sk{ग्रहः} (grahaḥ) — a planet}

\label{lesson:SA-C226-grahah}



\begin{quote}
Before the new one: say the Sanskrit for a bower, then the Sanskrit for a meadow.
\end{quote}

\subsection*{You'll want to know: \sk{ग्रहः}}
\textbf{\sk{ग्रहः}} — \emph{grahaḥ} — \textquotedblleft{}a planet\textquotedblright{}.

\begin{grammarlens}[title={What you've built}]
One more word for this chapter.
\end{grammarlens}

\subsection*{Your turn}
\begin{itemize}
\raggedright
  \item \emph{Say it:} \emph{grahaḥ}
  \item \emph{Say it:} \emph{grahaḥ}, once more
  \item \emph{Say it:} say \emph{grahaḥ}
  \item \emph{Recall:} say the Sanskrit for a bower, then the Sanskrit for a meadow, then say \emph{grahaḥ} again
\end{itemize}

\begin{tcolorbox}[breakable,colback=teal!4,colframe=teal!35!black,title={Before you move on}]
What does \sk{ग्रहः} mean? (\textquotedblleft{}A planet\textquotedblright{}.) Say it once more.
\end{tcolorbox}
\section[\texorpdfstring{\sk{धूमकेतुः} (dhūmaketuḥ)}{धूमकेतुः (dhūmaketuḥ)}]{\sk{धूमकेतुः} (dhūmaketuḥ) — a comet}

\label{lesson:SA-C226-dhumaketuh}



\begin{quote}
Before the new one: say the Sanskrit for a meadow, then the Sanskrit for a planet.
\end{quote}

\subsection*{You'll want to know: \sk{धूमकेतुः}}
\textbf{\sk{धूमकेतुः}} — \emph{dhūmaketuḥ} — \textquotedblleft{}a comet\textquotedblright{}.

\begin{grammarlens}[title={What you've built}]
One more word for this chapter.
\end{grammarlens}

\subsection*{Your turn}
\begin{itemize}
\raggedright
  \item \emph{Say it:} \emph{dhūmaketuḥ}
  \item \emph{Say it:} \emph{dhūmaketuḥ}, once more
  \item \emph{Say it:} say \emph{dhūmaketuḥ}
  \item \emph{Recall:} say the Sanskrit for a meadow, then the Sanskrit for a planet, then say \emph{dhūmaketuḥ} again
\end{itemize}

\begin{tcolorbox}[breakable,colback=teal!4,colframe=teal!35!black,title={Before you move on}]
What does \sk{धूमकेतुः} mean? (\textquotedblleft{}A comet\textquotedblright{}.) Say it once more.
\end{tcolorbox}
\section[\texorpdfstring{\sk{पशुः} (paśuḥ)}{पशुः (paśuḥ)}]{\sk{पशुः} (paśuḥ) — an animal}

\label{lesson:SA-C226-pasuh}



\begin{quote}
Before the new one: say the Sanskrit for a planet, then the Sanskrit for a comet.
\end{quote}

\subsection*{You'll want to know: \sk{पशुः}}
\textbf{\sk{पशुः}} — \emph{paśuḥ} — \textquotedblleft{}an animal\textquotedblright{}.

\begin{grammarlens}[title={What you've built}]
One more word for this chapter.
\end{grammarlens}

\subsection*{Your turn}
\begin{itemize}
\raggedright
  \item \emph{Say it:} \emph{paśuḥ}
  \item \emph{Say it:} \emph{paśuḥ}, once more
  \item \emph{Say it:} say \emph{paśuḥ}
  \item \emph{Recall:} say the Sanskrit for a planet, then the Sanskrit for a comet, then say \emph{paśuḥ} again
\end{itemize}

\begin{tcolorbox}[breakable,colback=teal!4,colframe=teal!35!black,title={Before you move on}]
What does \sk{पशुः} mean? (\textquotedblleft{}An animal\textquotedblright{}.) Say it once more.
\end{tcolorbox}
\section[\texorpdfstring{\sk{जन्तुः} (jantuḥ)}{जन्तुः (jantuḥ)}]{\sk{जन्तुः} (jantuḥ) — a creature}

\label{lesson:SA-C226-jantuh}



\begin{quote}
Before the new one: say the Sanskrit for a comet, then the Sanskrit for an animal.
\end{quote}

\subsection*{You'll want to know: \sk{जन्तुः}}
\textbf{\sk{जन्तुः}} — \emph{jantuḥ} — \textquotedblleft{}a creature\textquotedblright{}.

\begin{grammarlens}[title={What you've built}]
One more word for this chapter.
\end{grammarlens}

\subsection*{Your turn}
\begin{itemize}
\raggedright
  \item \emph{Say it:} \emph{jantuḥ}
  \item \emph{Say it:} \emph{jantuḥ}, once more
  \item \emph{Say it:} say \emph{jantuḥ}
  \item \emph{Recall:} say the Sanskrit for a comet, then the Sanskrit for an animal, then say \emph{jantuḥ} again
\end{itemize}

\begin{tcolorbox}[breakable,colback=teal!4,colframe=teal!35!black,title={Before you move on}]
What does \sk{जन्तुः} mean? (\textquotedblleft{}A creature\textquotedblright{}.) Say it once more.
\end{tcolorbox}
\section[\texorpdfstring{\sk{पक्षी} (pakṣī)}{पक्षी (pakṣī)}]{\sk{पक्षी} (pakṣī) — a winged creature, a bird}

\label{lesson:SA-C226-paksi}



\begin{quote}
Before the new one: say the Sanskrit for an animal, then the Sanskrit for a creature.
\end{quote}

\subsection*{You'll want to know: \sk{पक्षी}}
\textbf{\sk{पक्षी}} — \emph{pakṣī} — \textquotedblleft{}a winged creature, a bird\textquotedblright{}.

\begin{grammarlens}[title={What you've built}]
One more word for this chapter.
\end{grammarlens}

\subsection*{Your turn}
\begin{itemize}
\raggedright
  \item \emph{Say it:} \emph{pakṣī}
  \item \emph{Say it:} \emph{pakṣī}, once more
  \item \emph{Say it:} say \emph{pakṣī}
  \item \emph{Recall:} say the Sanskrit for an animal, then the Sanskrit for a creature, then say \emph{pakṣī} again
\end{itemize}

\begin{tcolorbox}[breakable,colback=teal!4,colframe=teal!35!black,title={Before you move on}]
What does \sk{पक्षी} mean? (\textquotedblleft{}A winged creature, a bird\textquotedblright{}.) Say it once more.
\end{tcolorbox}
\section[(dialogue)]{Words from chapters 221-223 — a first pass}

\label{lesson:SA-R226-first-pass-words-from-chapters-221-223}



\begin{quote}
Say the words for a creature and a winged creature.
\end{quote}

\subsection*{Your turn}
Say these aloud:

\begin{itemize}
\raggedright
  \item cotton, silk, leather, a knot, a small stool
  \item an instrument, a load, goods for sale, a shopkeeper, a newspaper
  \item a traveller, a guide, a holy ford, a hermitage, a monastery
\end{itemize}

\begin{tcolorbox}[breakable,colback=teal!4,colframe=teal!35!black,title={Before you move on}]
Cotton? (\textbf{\sk{कार्पासः}}, \emph{kārpāsaḥ}.) Silk? (\textbf{\sk{कौशेयम्}}, \emph{kauśeyam}.) Leather? (\textbf{\sk{चर्म}}, \emph{carma}.) A knot? (\textbf{\sk{ग्रन्थिः}}, \emph{granthiḥ}.) A small stool? (\textbf{\sk{पीठिका}}, \emph{pīṭhikā}.) An instrument? (\textbf{\sk{साधनम्}}, \emph{sādhanam}.) A load? (\textbf{\sk{भारः}}, \emph{bhāraḥ}.) Goods for sale? (\textbf{\sk{पण्यम्}}, \emph{paṇyam}.) A shopkeeper? (\textbf{\sk{आपणिकः}}, \emph{āpaṇikaḥ}.) A newspaper? (\textbf{\sk{पत्रिका}}, \emph{patrikā}.) A traveller? (\textbf{\sk{यात्रिकः}}, \emph{yātrikaḥ}.) A guide? (\textbf{\sk{मार्गदर्शकः}}, \emph{mārgadarśakaḥ}.) A holy ford? (\textbf{\sk{तीर्थम्}}, \emph{tīrtham}.) A hermitage? (\textbf{\sk{आश्रमः}}, \emph{āśramaḥ}.) A monastery? (\textbf{\sk{मठः}}, \emph{maṭhaḥ}.)
\end{tcolorbox}
\section[(dialogue)]{Words from chapters 224-226 — a second pass}

\label{lesson:SA-R226-second-pass-words-from-chapters-224-226}



\begin{quote}
Say the words for a centre and a tall building.
\end{quote}

\subsection*{Your turn}
Say these aloud:

\begin{itemize}
\raggedright
  \item a centre, a tall building, darkness, radiance, the ground
  \item a valley, a desert, a woodland, a bower, a meadow
  \item a planet, a comet, an animal, a creature, a winged creature
\end{itemize}

\begin{tcolorbox}[breakable,colback=teal!4,colframe=teal!35!black,title={Before you move on}]
A centre? (\textbf{\sk{केन्द्रम्}}, \emph{kendram}.) A tall building? (\textbf{\sk{अट्टालिका}}, \emph{aṭṭālikā}.) Darkness? (\textbf{\sk{अन्धकारः}}, \emph{andhakāraḥ}.) Radiance? (\textbf{\sk{ज्योतिः}}, \emph{jyotiḥ}.) The ground? (\textbf{\sk{भूमिः}}, \emph{bhūmiḥ}.) A valley? (\textbf{\sk{उपत्यका}}, \emph{upatyakā}.) A desert? (\textbf{\sk{मरुभूमिः}}, \emph{marubhūmiḥ}.) A woodland? (\textbf{\sk{काननम्}}, \emph{kānanam}.) A bower? (\textbf{\sk{कुञ्जः}}, \emph{kuñjaḥ}.) A meadow? (\textbf{\sk{शाद्वलम्}}, \emph{śādvalam}.) A planet? (\textbf{\sk{ग्रहः}}, \emph{grahaḥ}.) A comet? (\textbf{\sk{धूमकेतुः}}, \emph{dhūmaketuḥ}.) An animal? (\textbf{\sk{पशुः}}, \emph{paśuḥ}.) A creature? (\textbf{\sk{जन्तुः}}, \emph{jantuḥ}.) A winged creature? (\textbf{\sk{पक्षी}}, \emph{pakṣī}.)
\end{tcolorbox}
